\section{Introduction}

\subsection{Background and Objectives}

A search tree represents an ordered set while supporting insertion, deletion, and membership queries. The same set can have very different update costs because its representation depends on balance rules and access order. This project compares an unbalanced binary search tree (BST), an AVL tree, a splay tree, a red-black tree, and a B+ tree using the three update sequences required by Project 8-1, Compare Search Trees.

Each tree is implemented in its own C file, src/\allowbreak{}bst.c, src/\allowbreak{}avl.c, src/\allowbreak{}splay.c, src/\allowbreak{}redblack.c, and src/\allowbreak{}bplus.c, written in the style of the supplied testavl.c. Each file holds one node type and its operations. The driver in src/\allowbreak{}test.c includes those five files and runs either the smoke test or the timed benchmark.

\subsection{Problem Statement}

For each size N, the keys are the N distinct integers from 0 through N-1. The experiment inserts every key into an empty tree and then deletes every key. Scenario 1 inserts in increasing order and deletes in increasing order; Scenario 2 inserts in increasing order and deletes in decreasing order; Scenario 3 uses separately shuffled insertion and deletion permutations. The measured sizes are 1,000, 2,000, 5,000, 10,000, 20,000, 50,000, and 100,000.
